% sec_appendix.tex -- compact appendices for the main manuscript
\section{Strict recent-JFA comparison and background comparators}%
\label{app:related-work-scope}

This appendix separates the recent-\emph{J. Funct. Anal.} positioning from
background context.  The central present functional in the coefficient rows
is the fixed-law matched Poisson quotient
\[
  D_{\rm KL}(P_t*\nu\|P_t(\cdot-\bar{\gamma})).
\]
The comparison is about asymptotic coefficients of this matched-Poisson KL
quantity as \(t\to\infty\), not functional inequalities for the smoothed
measure \(P_t*\nu\).  The reproducibility convention for the recent-JFA
check is
\[
\begin{gathered}
\text{journal}=\text{JFA},\qquad
\text{years}=2021\text{--}2026,\\
\text{keywords}=\{\text{entropy, Poisson, Cauchy, KL, moment, quotient,}\\
\text{Bregman, transport}\}.
\end{gathered}
\]
The conclusion of that check is theorem-level rather than bibliographic: the
nearest recent-JFA papers address changing-law limit theorems, linear heat
profiles, Hardy--Stein identities, transport maps, or functional
inequalities, but not continuous Cauchy-kernel KL coefficient asymptotics for
one fixed input law.  Older JFA papers and non-JFA papers remain useful
background, and are recorded separately in
Table~\ref{tab:older-extra-background}.

\subsection*{JFA 2021--2026 theorem-level comparison}

{\scriptsize
\setlength{\tabcolsep}{2pt}
\renewcommand{\arraystretch}{1.22}
\begin{longtable}{>{\raggedright\arraybackslash}p{0.20\linewidth}|%
>{\raggedright\arraybackslash}p{0.19\linewidth}|%
>{\raggedright\arraybackslash}p{0.28\linewidth}|%
>{\raggedright\arraybackslash}p{0.27\linewidth}}
\caption{Recent \emph{J. Funct. Anal.} theorem-level comparison.  Every row
uses the protocol
\(\text{journal}=\text{JFA}\), \(\text{years}=2021\text{--}2026\), and
\(\text{keywords}=\{\text{entropy, Poisson, Cauchy, KL, moment, quotient,
Bregman, transport}\}\); exactly one closest checked recent-JFA comparator is
listed for each present theorem package.}
\label{tab:recent-jfa-comparison}\\
\textbf{present theorem package} & \textbf{protocol}
& \textbf{closest checked recent-JFA comparator}
& \textbf{one-sentence non-overlap}\\
\hline
Theorem~I and Theorem~\ref{thm:poisson-kl-even-orders}, including the
finite-variance threshold.
& JFA; 2021--2026; keywords
\(\{\)entropy, Poisson, Cauchy, KL, moment, quotient, Bregman, transport\(\}\).
& Bobkov--G{\"o}tze, \emph{Esscher transform and the central limit theorem}
\cite{BobkovGotzeEsscher2025};
\nolinkurl{doi:10.1016/j.jfa.2025.110999},
\nolinkurl{arXiv:2407.20726}.
& It studies changing normalized sums and Esscher transforms, whereas
Theorem~I fixes \(\nu\), varies the Poisson scale \(t\), computes
Cayley--Chebyshev Laurent coefficients, and proves the fixed-law leading
finite-KL threshold.\\
\hline
Transfer Lemma~I.1 and Corollaries~I.2--I.4 on verified missing-top-moment
tail replacement.
& JFA; 2021--2026; keywords
\(\{\)entropy, Poisson, Cauchy, KL, moment, quotient, Bregman, transport\(\}\).
& Bobkov--G{\"o}tze, \emph{Esscher transform and the central limit theorem}
\cite{BobkovGotzeEsscher2025};
\nolinkurl{doi:10.1016/j.jfa.2025.110999},
\nolinkurl{arXiv:2407.20726}.
& It has no fixed-law quotient input of the form
\(\Delta_t=A_t+u_pt^{-p}M_p(t)+e_t\) with
\(\|e_t\|_\infty+\|e_t\|_{L^1(\omega)}=o(t^{-p}M_p(t))\), which is the
mechanism used here for exact, perturbative, and monotone slowly varying
tails.\\
\hline
Theorem~II and Theorem~\ref{thm:multidimensional-poisson-entropy-jets}.
& JFA; 2021--2026; keywords
\(\{\)entropy, Poisson, Cauchy, KL, moment, quotient, Bregman, transport\(\}\).
& Anker--Papageorgiou--Zhang
\cite{AnkerPapageorgiouZhang2023HeatSymmetricSpaces};
\nolinkurl{doi:10.1016/j.jfa.2022.109828},
\nolinkurl{arXiv:2112.01323}.
& It treats linear heat profiles on noncompact symmetric spaces, not a
Euclidean radial matched-Poisson quotient jet converted into nonlinear tensor
KL coefficients \(\sum_{m=4}^{L+2}\mathcal A_m^{(d)}t^{-m}\).\\
\hline
Corollary~II.1 and Theorem~\ref{thm:multidimensional-density-quotient}.
& JFA; 2021--2026; keywords
\(\{\)entropy, Poisson, Cauchy, KL, moment, quotient, Bregman, transport\(\}\).
& Anker--Papageorgiou--Zhang
\cite{AnkerPapageorgiouZhang2023HeatSymmetricSpaces};
\nolinkurl{doi:10.1016/j.jfa.2022.109828},
\nolinkurl{arXiv:2112.01323}.
& It does not prove that finite centred \(L\)-absolute moment in \(\RR^d\),
with \(L\ge d+1\), supplies the matched radial Poisson quotient jet in
\(L^\infty\cap L^1(\Omega_d)\).\\
\hline
Proposition~II.2 and
Proposition~\ref{prop:multidim-formal-covariance-layer}.
& JFA; 2021--2026; keywords
\(\{\)entropy, Poisson, Cauchy, KL, moment, quotient, Bregman, transport\(\}\).
& Cheng--Thalmaier--Wang, \emph{Entropy, Fisher information and Stein
discrepancy on Riemannian manifolds}
\cite{ChengThalmaierWang2023EntropySteinTransport};
\nolinkurl{doi:10.1016/j.jfa.2023.109997},
\nolinkurl{arXiv:2108.12755}.
& It proves entropy/Fisher/Stein/transport inequalities, not the explicit
radial-Poisson covariance contraction
\(\mathcal Q_d(\Sigma)=c_d^{\rm iso}(\operatorname{tr}\Sigma)^2+
c_d^{\rm tr}\|\Sigma_0\|_{\rm HS}^2\), and finite covariance alone is not a
KL theorem here.\\
\hline
Corollary~II.3 and
Corollary~\ref{cor:finite-dplusone-moment-leading-multidim-coefficient}.
& JFA; 2021--2026; keywords
\(\{\)entropy, Poisson, Cauchy, KL, moment, quotient, Bregman, transport\(\}\).
& Chen--Chewi--Niles-Weed
\cite{ChenChewiNilesWeed2021MixtureLogSobolev};
\nolinkurl{doi:10.1016/j.jfa.2021.109236},
\nolinkurl{arXiv:2102.11476}.
& It proves dimension-free log-Sobolev inequalities for mixtures, whereas the
present corollary uses finite \((d+1)\)-moment only to verify the quotient
criterion before identifying the \(t^{-4}\) matched-KL coefficient.\\
\hline
Theorem~III and Theorem~\ref{thm:poisson-kl-doob-tail}.
& JFA; 2021--2026; keywords
\(\{\)entropy, Poisson, Cauchy, KL, moment, quotient, Bregman, transport\(\}\).
& Bogdan--Gutowski--Pietruska-Pa{\l}uba
\cite{BogdanGutowskiPietruskaPaluba2025PolarizedHardyStein};
\nolinkurl{doi:10.1016/j.jfa.2025.110827},
\nolinkurl{arXiv:2309.09856}.
& It supplies a polarized Hardy--Stein identity, not the compact-window
verification of the moving matched quotient
\((P_t*\nu)/P_t(\cdot-\bar{\gamma})\) or the exact Cauchy-scale KL tail
formula \(D_{\rm KL}(h_t\|g_t)=\int_t^\infty\mathcal I_\nu(s)\dd s\).\\
\hline
Corollary~I.5 and Corollary~\ref{cor:coefficient-defect-coupling}.
& JFA; 2021--2026; keywords
\(\{\)entropy, Poisson, Cauchy, KL, moment, quotient, Bregman, transport\(\}\).
& Cheng--Thalmaier--Wang, \emph{Entropy, Fisher information and Stein
discrepancy on Riemannian manifolds}
\cite{ChengThalmaierWang2023EntropySteinTransport};
\nolinkurl{doi:10.1016/j.jfa.2023.109997},
\nolinkurl{arXiv:2108.12755}.
& It proves Stein and transport inequalities, not the explicit \(A_6\)
Laurent coefficient or the resulting non-sharp finite-moment coupling to the
symmetric two-point law.\\
\hline
Continuous Cauchy-kernel uses of the word ``Poisson'' throughout the theorem
package.
& JFA; 2021--2026; keywords
\(\{\)entropy, Poisson, Cauchy, KL, moment, quotient, Bregman, transport\(\}\).
& L{\'o}pez-Rivera--Shenfeld, \emph{The Poisson transport map}
\cite{LopezRiveraShenfeld2025PoissonTransport};
\nolinkurl{doi:10.1016/j.jfa.2025.110864},
\nolinkurl{arXiv:2407.02359}.
& It concerns contraction transport from Poisson point processes to
ultra-log-concave measures on \(\mathbb N\), while the present ``Poisson''
object is the continuous real-line Cauchy/Poisson kernel \(P_t\) and its
matched KL coefficient asymptotics.\\
\end{longtable}
}

\subsection*{Older background and non-JFA comparators}

{\scriptsize
\setlength{\tabcolsep}{2pt}
\renewcommand{\arraystretch}{1.2}
\begin{longtable}{>{\raggedright\arraybackslash}p{0.18\linewidth}|%
>{\raggedright\arraybackslash}p{0.24\linewidth}|%
>{\raggedright\arraybackslash}p{0.24\linewidth}|%
>{\raggedright\arraybackslash}p{0.24\linewidth}}
\caption{Older background and non-JFA comparators.  These references clarify
language, motivation, or proof technology, but they are not part of the
strict 2021--2026 JFA theorem-level comparison above.}
\label{tab:older-extra-background}\\
\textbf{present point} & \textbf{background reference} &
\textbf{use in the manuscript} & \textbf{scope separation}\\
\hline
Corollary~I.5 and
Corollary~\ref{cor:coefficient-defect-coupling}.
& Nourdin--Peccati--Swan, \emph{Entropy and the fourth moment phenomenon}
\cite{NourdinPeccatiSwan2014EntropyFourthMoment}.
& Supplies older JFA context for the phrase ``fourth-moment defect'' and for
entropy phenomena tied to fourth moments.
& It is a 2014 JFA paper, so it is deliberately kept out of the
2021--2026 comparison.  Its theorem class is also different: Gaussian
fourth-moment phenomena on Wiener chaos, not fixed-law Poisson quotient
Laurent coefficients or the \(A_6,A_8\) matched-Cauchy calculation.\\
\hline
Theorem~III and the Bregman tail identity.
& Hardy--Stein, Beurling--Deny, and stable de Bruijn references such as
\cite{Gutowski2023HardySteinPureJump,
GutowskiKwasnicki2025BeurlingDenySobolevBregman,
JohnsonStableDeBruijn2013}, together with the information-estimation
de Bruijn line
\cite{GuoShamaiVerdu2005MMSE,Verdu2010Mismatched}.
& Provide proof language for nonlocal Bregman forms, cutoff identities, and
entropy dissipation.  The Gaussian-channel identities are included only as
background for differentiating information along smoothing/noise flows.
& They are not theorem-level competitors for the moving matched quotient
\(u_t=(P_t*\nu)/P_t(\cdot-\bar{\gamma})\), whose compact-window bounds and
tail identity are proved here; no MMSE or Gaussian-channel estimation formula
is used in the coefficient proofs.\\
\hline
Gaussian-convolution and Bakry--{\'E}mery log-Sobolev context.
& Gaussian-convolution log-Sobolev and diffusion-curvature references such as
\cite{BardetGozlanMalrieuZitt2018GaussianConvolutions,
BakryEmery1985Diffusions,BakryGentilLedoux2014}.
& Provide comparison language for functional inequalities, curvature
stability, and the difference between smoothing inequalities and asymptotic
coefficient extraction.
& These references concern Gaussian convolutions or Markov diffusion
semigroups.  The present paper studies matched Cauchy/Poisson quotient
coefficients for one fixed input law and does not derive its results from
Bakry--{\'E}mery stability or a Gaussian-convolution log-Sobolev theorem.\\
\hline
The finite-moment and power-tail quotient expansions.
& Euclidean heat, fractional-heat, Poisson-kernel, and regular-variation
sources
\cite{EscobedoZuazua1991,DuoandikoetxeaZuazua1992,
BilerKarchWoyczynski2001Levy,Yamamoto2012FractionalExpansion,
IshigeKawakamiMichihisa2017,IshigeKawakamiMichihisa2022RefinedFractional,
FilaIshigeKawakami2012PoissonKernel,Iwabuchi2015CriticalBurgersPoisson,
Papageorgiou2024FractionalLargeTime,BinghamGoldieTeugels1987,Resnick1987}.
& Motivate profile-expansion and truncated-moment terminology used before
the nonlinear KL transfer.
& They do not replace the explicit matched-Poisson quotient remainders,
finite Laurent arithmetic, or restricted slowly varying hypotheses stated in
the present theorems.\\
\hline
Cauchy-information normalizations.
& Cauchy KL and \(f\)-divergence sources
\cite{ChyzakNielsen2019CauchyKL,NielsenOkamura2023Cauchy,Verdu2023Cauchy}.
& Fix static Cauchy-parameter context and information-geometric notation.
& The coefficient hierarchy here concerns the non-Cauchy mixture quotient
\(P_t*\nu\) against its matched translated Poisson profile, not divergences
between two Cauchy laws.\\
\end{longtable}
}

\subsection*{Scope restrictions}

The consolidated scope ledger is Subsection~\ref{subsec:intro-hypotheses-scope}.
The appendix does not restate it.  The tables above are comparison and
background records; they should be read together with that subsection for the
conditional quotient hypotheses, finite-covariance criterion, tail-subclass
restrictions, and non-sharp coefficient-defect interpretation used by the
main theorem package.

\subsection*{Background used in proofs}

Older and non-JFA references are used as proof language rather than as
theorem-level competitors.  Cauchy KL and \(f\)-divergence formulae
\cite{ChyzakNielsen2019CauchyKL,NielsenOkamura2023Cauchy,Verdu2023Cauchy}
give static Cauchy-parameter context, whereas
Theorem~\ref{thm:poisson-kl-even-orders} expands the
non-Cauchy mixture quotient \(P_t*\nu\) against its matched translated
Poisson profile.  Gaussian-convolution log-Sobolev results
\cite{BardetGozlanMalrieuZitt2018GaussianConvolutions,
ChenChewiNilesWeed2021MixtureLogSobolev} and the Bakry--{\'E}mery diffusion
framework \cite{BakryEmery1985Diffusions,BakryGentilLedoux2014} are cited
only to locate neighbouring functional-inequality questions; they do not
provide the quotient expansion or the Laurent coefficient arithmetic.
Likewise, de Bruijn and information-estimation identities
\cite{DeBruijn1953,GuoShamaiVerdu2005MMSE,Verdu2010Mismatched} give
background for entropy differentiation along smoothing channels, while the
Cauchy Bregman identity here is proved directly from the Poisson generator,
the compact-window quotient bounds, and the explicit \(L^1\) cutoff
majorant.  Euclidean heat, fractional-heat, and Poisson-kernel moment
expansions
\cite{EscobedoZuazua1991,DuoandikoetxeaZuazua1992,
BilerKarchWoyczynski2001Levy,Yamamoto2012FractionalExpansion,
IshigeKawakamiMichihisa2017,IshigeKawakamiMichihisa2022RefinedFractional,
FilaIshigeKawakami2012PoissonKernel,Iwabuchi2015CriticalBurgersPoisson,
Papageorgiou2024FractionalLargeTime}
motivate finite quotient jets before the nonlinear KL transfer.  Regular
variation references
\cite{Feller1971,BinghamGoldieTeugels1987,Resnick1987,Resnick2007HeavyTail}
fix the language of truncated moments; the paper proves the exact KL variance
threshold, the displayed power-tail cases, and one conditional slowly varying
add-on only under explicit quotient-remainder hypotheses.  Entropic
stable-CLT and Wasserstein/Mallows limit results
\cite{BobkovChistyakovGotze2013,BobkovChistyakovGotzeEdgeworth2013,
BobkovChistyakovGotzeBerryEsseen2014,BobkovGotzeRenyiSurvey2025,
BobkovGotzeRenyi2025,BobkovChistyakovGotzeStableFisher2014,
JohnsonStableDeBruijn2013,Toscani2016,Mallows1972AsymptoticJointNormality,
BickelFreedman1981Bootstrap,JohnsonSamworth2005MallowsStable,
BarbosaDorea2009MallowsStable,DoreaFerreira2012MallowsStable,
SamorodnitskyTaqqu1994Stable}
concern changing laws; here the law is fixed and only the Poisson scale
changes.  De Bruijn--Stam, fractional-generator, and Hardy--Stein/Bregman
sources
\cite{DeBruijn1953,Stam1959,Blachman1965,Barron1986,JohnsonBarron2004,
MadimanBarron2007,Verdu2010Mismatched,
HirataNemotoYoshida2012RelativeEntropy,Yoshida2017RelativeEntropyDiffusion,
CaffarelliSilvestre2007,StingaTorrea2010,Kwasnicki2017,
BanuelosBogdanLuks2016HardyStein,BanuelosKim2019HardySteinNonsymmetric,
BogdanGrzywnyPietruskaPalubaRutkowski2023NonlinearNonlocalDouglas,
Gutowski2023HardySteinPureJump,GutowskiKwasnicki2025BeurlingDenySobolevBregman}
provide the dissipation vocabulary; the new step here is the direct
annular-cutoff domain verification for the matched Cauchy quotient, followed
by the exact tail identity and coefficient transfer.

\section{Weighted trigonometric normalization}%
\label{app:entropy-integrals}

This appendix records the change of variables behind the Laurent coefficient
normalization.  Set \(y=\tan\theta\),
\(-\pi/2<\theta<\pi/2\).  Then
\[
  \frac{\dd y}{\pi(1+y^2)}=\frac{\dd\theta}{\pi}.
\]
With \(z=e^{2\ii\theta}\),
\[
  \cos(2k\theta)=\frac{z^k+z^{-k}}2,\qquad
  \sin(2k\theta)=\frac{z^k-z^{-k}}{2\ii}.
\]
Thus the Cauchy mode integrals in
Theorem~\ref{thm:finite-fourier-mode-integral-algorithm} are exactly Haar
constant terms:
\[
  \int_{\RR}\prod_j u_{n_j}(y)\,\omega(\dd y)
  =[z^0]\prod_j Q_{n_j}(z).
\]
For order eight, Lemma~\ref{lem:eighth-order-coefficient-arithmetic} gives
the recurrence; Appendix~\ref{app:order-eight-laurent-certificate} gives the
row certificate.

\section{Order-eight Laurent certificate}%
\label{app:order-eight-laurent-certificate}

This appendix gives the finite certificate behind
Lemma~\ref{lem:eighth-order-coefficient-arithmetic}.  For
\(\mathbf n=(n_1,\ldots,n_r)\), put
\[
\mathcal K(\mathbf n)
:=\{(k_1,\ldots,k_r):0<|k_j|\le n_j,\ k_1+\cdots+k_r=0\}.
\]
With \(Q_n(z)=\sum_kq_{n,k}z^k\) and the uniform coefficients
\eqref{eq:mode-laurent-coefficients}, put
\[
\lambda_n:=2^{-n}(-\ii)^n .
\]
For each row set
\[
S(\mathbf n):=
\sum_{\mathbf k\in\mathcal K(\mathbf n)}
\prod_{j=1}^r
\operatorname{sgn}(k_j)^{\,n_j\bmod2}
\binom{n_j-1}{|k_j|-1},
\qquad
\Lambda(\mathbf n):=\prod_{j=1}^r\lambda_{n_j}.
\]
Then the all-sign coefficient convention gives
\[
J(\mathbf n)=[z^0]\prod_{j=1}^rQ_{n_j}(z)
=\Lambda(\mathbf n)S(\mathbf n).
\]

The Laurent representatives used in every row of the \(A_6,A_8\)
substitution are:
\begin{align}
Q_2(z)
&=-\frac14\bigl(z+z^{-1}+z^2+z^{-2}\bigr),
\label{eq:appendix-d-explicit-q2}\\
Q_3(z)
&=\frac{\ii}{8}\bigl((z-z^{-1})+2(z^2-z^{-2})+(z^3-z^{-3})\bigr),
\label{eq:appendix-d-explicit-q3}\\
Q_4(z)
&=\frac1{16}\bigl((z+z^{-1})+3(z^2+z^{-2})
  +3(z^3+z^{-3})+(z^4+z^{-4})\bigr),
\label{eq:appendix-d-explicit-q4}\\
Q_5(z)
&=-\frac{\ii}{32}\bigl((z-z^{-1})+4(z^2-z^{-2})
  +6(z^3-z^{-3})+4(z^4-z^{-4})+(z^5-z^{-5})\bigr),
\label{eq:appendix-d-explicit-q5}\\
Q_6(z)
&=-\frac1{64}\bigl((z+z^{-1})+5(z^2+z^{-2})
  +10(z^3+z^{-3})+10(z^4+z^{-4}) \nonumber\\
&\hspace{6.7em}
  +5(z^5+z^{-5})+(z^6+z^{-6})\bigr).
\label{eq:appendix-d-explicit-q6}
\end{align}
These are exactly the cases \(2\le n\le6\) of
\eqref{eq:mode-laurent-polynomial-definition}.  The display is included here
so that the signs in the odd negative-exponent rows can be checked directly
against the table below.

Here is an independent checksum for the displayed Laurent representatives and
for three rows used later in the certificate.  For a Laurent polynomial
\(F(z)=\sum_{|k|\le K}a_kz^k\), put
\[
\mathsf G_N(F):=\frac1N\sum_{\ell=0}^{N-1}F(\zeta_N^\ell),
\qquad \zeta_N:=e^{2\pi\ii/N}.
\]
If \(N>K\), then \(\mathsf G_N(F)=[z^0]F\), because no nonzero exponent in
\([-K,K]\) is divisible by \(N\).  Thus the following root-of-unity checks
extract the same constant terms without using the zero-sum recurrence below:
\[
\begin{array}{c|c|c}
\text{object}&\text{checksum}&\text{value}\\
\hline
Q_2&\bigl(Q_2(1),Q_2(-1)\bigr)&(-1,0)\\
Q_3&\bigl(Q_3(1),Q_3(-1)\bigr)&(0,0)\\
Q_4&\bigl(Q_4(1),Q_4(-1)\bigr)&(1,0)\\
Q_5&\bigl(Q_5(1),Q_5(-1)\bigr)&(0,0)\\
Q_6&\bigl(Q_6(1),Q_6(-1)\bigr)&(-1,0)\\
J(2,4)&\mathsf G_{13}(Q_2Q_4)&-1/8\\
J(3,5)&\mathsf G_{17}(Q_3Q_5)&-15/128\\
J(2,3,3)&\mathsf G_{17}(Q_2Q_3^2)&-9/128 .
\end{array}
\]
The endpoint checks verify the parity and total-sign normalization of
\(Q_2,\ldots,Q_6\); the grid checks use \(13>2+4\) and
\(17>3+5=2+3+3\), so they are exact finite averages rather than numerical
quadrature.

\begin{lemma}[Deterministic exact Laurent recurrence]%
\label{lem:appendix-d-deterministic-laurent-recurrence}
Let \(\mathbf n=(n_1,\ldots,n_r)\) be any finite row of positive integers.
Define finitely supported maps \(C_j:\ZZ\to\QQ[\ii]\) by
\[
C_0(s)=\mathbf 1_{\{s=0\}},\qquad
C_j(s)=\sum_{0<|k|\le n_j}
2^{-n_j}(-\ii)^{n_j}\operatorname{sgn}(k)^{\,n_j\bmod2}
\binom{n_j-1}{|k|-1}\,C_{j-1}(s-k).
\]
Then
\[
  C_r(0)=[z^0]\prod_{j=1}^r Q_{n_j}(z)=J(\mathbf n).
\]
In particular \(C_r(0)\in\QQ\).  More explicitly,
\[
C_r(0)=
\prod_{j=1}^r\lambda_{n_j}
\sum_{\mathbf k\in\mathcal K(\mathbf n)}
\prod_{j=1}^r
\operatorname{sgn}(k_j)^{\,n_j\bmod2}
\binom{n_j-1}{|k_j|-1}.
\]
The recurrence is deterministic and terminates after at most
\(\prod_{j=1}^r(2n_j)\) rational Gaussian-integer operations.
\end{lemma}

\begin{proof}
The coefficient formula \eqref{eq:mode-laurent-coefficients} says that
multiplication by \(Q_{n_j}\) sends a coefficient array \(C_{j-1}\) to the
convolution array displayed above.  Induction on \(j\) gives
\[
  C_j(s)=[z^s]\prod_{a=1}^jQ_{n_a}(z),
\]
and hence \(C_r(0)\) is the desired constant term.  Expanding the last
coefficient and collecting the common factor
\(\prod_j\lambda_{n_j}\) gives the stated zero-sum formula.

The value is rational because the constant term is the Haar integral of the
real-valued product
\(\prod_j u_{n_j}(\tan\theta)\), by
Theorem~\ref{thm:finite-fourier-mode-integral-algorithm}.  Equivalently,
when \(\sum_j n_j\) is odd, the integrand is odd and the value is zero; when
\(\sum_j n_j\) is even, \(\prod_j\lambda_{n_j}\) is real and the zero-sum
integer convolution is an integer.
\end{proof}

Thus the certificate used in the proof is completely finite: for each row
below, the integer \(S(\mathbf n)\) is obtained by the zero-sum convolution
in Lemma~\ref{lem:appendix-d-deterministic-laurent-recurrence}, multiplication
by the displayed rational factor \(\Lambda(\mathbf n)\) gives the exact
constant term \(J(\mathbf n)\), and no external computation is used in the
argument.

For the mode products entering \(A_4,A_6,A_8\), the finite zero-sum
convolution is:
\[
\begin{array}{c|r|r|c|c}
\mathbf n&|\mathcal K(\mathbf n)|&S(\mathbf n)&\Lambda(\mathbf n)&J(\mathbf n)\\
\hline
(2,2)&4&4&1/16&1/4\\
(2,4)&4&8&-1/64&-1/8\\
(3,3)&6&-12&-1/64&3/16\\
(2,2,2)&6&6&-1/64&-3/32\\
(2,6)&4&12&1/256&3/64\\
(3,5)&6&-30&1/256&-15/128\\
(4,4)&8&40&1/256&5/32\\
(2,2,4)&12&24&1/256&3/32\\
(2,3,3)&14&-18&1/256&-9/128\\
(2,2,2,2)&36&36&1/256&9/64 .
\end{array}
\]
Odd total-order entries through order eight vanish by
\(u_n(-y)=(-1)^nu_n(y)\).  The sign-sensitive rows
\((3,5)\) and \((2,3,3)\) use the same all-sign coefficient rule as the
other rows; no separate conjugation convention is applied.

For reproducibility it is enough to combine the recurrence with the row table.
The following short checks record the sign-sensitive convention used in the
rows that most easily go wrong.  Write
\[
b_m^{(n)}:=\binom{n-1}{m-1}\qquad(1\le m\le n).
\]
First, the lowest row gives
\[
S(2,2)=4,\qquad
J(2,2)=\Lambda(2,2)S(2,2)=\frac{4}{16}=\frac14,
\]
so \(A_4=\frac12\mu_2^2J(2,2)=\mu_2^2/8\).  For the odd two-factor row
\((3,5)\), the zero-sum condition forces \(\ell=-k\).  The two odd signs
therefore contribute
\(\operatorname{sgn}(k)\operatorname{sgn}(-k)=-1\), and
\[
S(3,5)
=-2\sum_{m=1}^{3}b_m^{(3)}b_m^{(5)}
=-2(1+8+6)=-30.
\]
Since \(\Lambda(3,5)=(\ii/8)(-\ii/32)=1/256\), this gives
\[
J(3,5)=\frac{-30}{256}=-\frac{15}{128}.
\]

For the mixed row \((2,3,3)\), set \(a\in\{\pm1,\pm2\}\) for the \(Q_2\)
index and sum over \(b+c=-a\) for the two \(Q_3\) indices.  The \(Q_2\)
factor has no sign, while the two odd factors contribute
\(\operatorname{sgn}(b)\operatorname{sgn}(c)\).  The four outer cases are
\[
\begin{array}{c|c}
a&
\displaystyle
\sum_{\substack{0<|b|,|c|\le3\\ b+c=-a}}
\operatorname{sgn}(b)\operatorname{sgn}(c)
b_{|b|}^{(3)}b_{|c|}^{(3)}\\
\hline
1&-8\\
-1&-8\\
2&-1\\
-2&-1 .
\end{array}
\]
Thus
\[
S(2,3,3)=-8-8-1-1=-18.
\]
Here \(\Lambda(2,3,3)=(-1/4)(\ii/8)^2=1/256\), so
\[
J(2,3,3)=\frac{-18}{256}=-\frac9{128}.
\]
All remaining rows are obtained in the same way by the deterministic
convolution of Lemma~\ref{lem:appendix-d-deterministic-laurent-recurrence}.
Thus the submitted certificate consists of the signed coefficient rule, the
finite recurrence, the checksum display, and the row table.  No unsubmitted
computation is invoked; only rational Gaussian-integer arithmetic is involved,
the imaginary parts cancel in the even total-order rows listed above, and odd
total-order rows vanish by parity.

\subsection*{Optional order-eight coefficient layer}
\label{subsec:appendix-optional-order-eight-layer}

The following statements record the order-eight consequence of the same row
table.  They are not used in the sixth-order defect coupling of
Subsection~\ref{subsec:coefficient-defect-coupling}; the \(A_6\) defect theorem
depends only on the rows of total order at most six.

\begin{lemma}[Coefficients through order eight from the row table]%
\label{lem:eighth-order-coefficient-arithmetic}
Applying Lemma~\ref{lem:appendix-d-deterministic-laurent-recurrence} to the
rows entering \(A_4,A_6,A_8\) gives the table displayed above, in particular
\[
J(2,2)=\frac14,\qquad
J(2,4)=-\frac18,\qquad
J(3,3)=\frac3{16},
\]
\[
J(2,3,3)=-\frac9{128},\qquad
J(3,5)=-\frac{15}{128}.
\]
Substitution gives
\[
A_4=\frac{\sigma^4}{8},\qquad
A_6=\frac{\sigma^6+6\mu_3^2-8\sigma^2\mu_4}{64},
\]
and
\[
A_8=
\frac{3\sigma^8-12\sigma^4\mu_4+9\sigma^2\mu_3^2
+12\sigma^2\mu_6-30\mu_3\mu_5+20\mu_4^2}{256}.
\]
\end{lemma}

\begin{proof}
The recurrence is exactly
Lemma~\ref{lem:appendix-d-deterministic-laurent-recurrence} specialized to
the rows of total order at most eight.  Combining those row values with the
universal coefficient formula \eqref{eq:universal-entropy-coefficient} gives
\[
\begin{aligned}
A_4&=\frac12\mu_2^2J(2,2),\\
A_6&=\mu_2\mu_4J(2,4)+\frac12\mu_3^2J(3,3)
-\frac16\mu_2^3J(2,2,2),\\
A_8&=\mu_2\mu_6J(2,6)+\mu_3\mu_5J(3,5)
+\frac12\mu_4^2J(4,4)
-\frac12\mu_2^2\mu_4J(2,2,4)\\
&\hspace{2.5em}
-\frac12\mu_2\mu_3^2J(2,3,3)
+\frac1{12}\mu_2^4J(2,2,2,2).
\end{aligned}
\]
The row table above gives the displayed rational polynomials.  The parity
identity \(u_n(-y)=(-1)^nu_n(y)\) eliminates all odd total-order products
through order eight.
\end{proof}

\begin{lemma}[Substitution into \(A_4,A_6,A_8\)]%
\label{lem:n4-entropy-coefficient-substitution}
For \(\mu_2=\sigma^2\),
\[
A_4=\frac{\sigma^4}{8},
\qquad
A_6=\frac{\sigma^6+6\mu_3^2-8\sigma^2\mu_4}{64},
\]
and
\[
A_8=
\frac{3\sigma^8-12\sigma^4\mu_4+9\sigma^2\mu_3^2
+12\sigma^2\mu_6-30\mu_3\mu_5+20\mu_4^2}{256}.
\]
\end{lemma}

\begin{proof}
This is the substitution just carried out in
Lemma~\ref{lem:eighth-order-coefficient-arithmetic}.
\end{proof}

\begin{theorem}[Optional coefficient identification through order eight]%
\label{thm:poisson-kl-eighth}
Let \(\nu\) have mean \(\bar{\gamma}\) and finite variance
\(\sigma^2\ge0\).  If \(\nu\) has finite centred fourth absolute moment, then
\[
D_{\rm KL}(P_t*\nu\|P_t(\cdot-\bar{\gamma}))
=\frac{\sigma^4}{8t^4}
+\frac{\sigma^6+6\mu_3^2-8\sigma^2\mu_4}{64t^6}
+o(t^{-6}).
\]
If \(\nu\) has finite centred sixth absolute moment, then
\begin{equation}\label{eq:poisson-kl-eighth}
D_{\rm KL}(P_t*\nu\|P_t(\cdot-\bar{\gamma}))
=\frac{\sigma^4}{8t^4}+\frac{A_6(\nu)}{t^6}
+\frac{A_8(\nu)}{t^8}+o(t^{-8}),
\end{equation}
where \(A_6,A_8\) are the polynomials in
Lemma~\ref{lem:n4-entropy-coefficient-substitution}.  When
\(\sigma^2=0\), these are the trivial exact identities for the point mass
\(\delta_{\bar{\gamma}}\).
\end{theorem}

\begin{proof}
The case \(\sigma^2=0\) is immediate.  Otherwise apply
Theorem~\ref{thm:poisson-kl-even-orders} with \(N=3\) for the sixth
coefficient and \(N=4\) for the eighth coefficient, then use
Lemma~\ref{lem:n4-entropy-coefficient-substitution}.
\end{proof}

\begin{theorem}[Optional \(N=4\) remainder under the seventh moment]%
\label{thm:poisson-kl-eighth-seventh-remainder}
If, in addition, the centred seventh absolute moment is finite, then
\[
D_{\rm KL}(P_t*\nu\|P_t(\cdot-\bar{\gamma}))
=\frac{\sigma^4}{8t^4}+\frac{A_6(\nu)}{t^6}
+\frac{A_8(\nu)}{t^8}+O_\nu(t^{-9}).
\]
\end{theorem}

\begin{proof}
Apply Lemma~\ref{lem:positive-margin-quotient-remainder} with
\(N=4\), \(L=6\), and \(\eta=1\).  The seventh absolute moment gives
\[
\delta_t(y)=\sum_{n=2}^{6}\mu_nu_n(y)t^{-n}+\rho_{6,t}(y),
\qquad
\|\rho_{6,t}\|_\infty+\|\rho_{6,t}\|_{L^1(\omega)}=O_\nu(t^{-7}),
\]
and \(\|\delta_t\|_\infty=O_\nu(t^{-2})\).  Every entropy term containing
\(\rho_{6,t}\) has one further density factor at the quadratic entropy scale,
so its contribution is \(O_\nu(t^{-9})\), while the Taylor tail after degree
\(4\) is \(O_\nu(t^{-10})\).  The coefficient theorem above supplies the
pure-mode terms \(A_4,A_6,A_8\).
\end{proof}
